\documentclass[conference]{IEEEtran}
\usepackage[utf8]{inputenc}
\usepackage[T1]{fontenc}
\usepackage{cite}
\usepackage{amsmath,amssymb}
\usepackage{graphicx}
\usepackage{booktabs}
\usepackage{url}

\title{Verifier-Gated LLM NetOps: A Preregistered, Artifact-Backed Evaluation of Input-Sanitization and Deterministic Verification Against Adversarial Intents}

\author{
\IEEEauthorblockN{Autonomous AI Researcher}
\IEEEauthorblockA{CNSM 2026 Agentic AI Researcher Track}
}

\begin{document}

\maketitle

\begin{abstract}
We report a preregistered paired confirmatory experiment (study\_id cand-2026-03) that measured whether template-based input sanitization combined with deterministic verifier-gating (and at-most-one automated repair) changes the rate at which a deterministic validator accepts LLM-generated CLI configurations under adversarially mutated intents. The frozen run used the declared generator gpt-5-mini and deterministic\_netops\_validator\_v5 within the hosted\_netops\_gvr\_v1 adapter. Planned episodes = 300; completed episodes = 282; complete paired outcomes used in the primary analysis = 141. Observed per-pair acceptance: Baseline 135/141 (95.7447\%); Guarded 141/141 (100.0\%); paired absolute difference = 0.04255319148936165; exact two-sided McNemar p = 0.03125; 95\% bootstrap percentile CI (10,000 resamples, seed = 1729) = [0.014184397163120567, 0.07801418439716312]. We reconcile preregistration language vs the formal estimand, document per-condition pipeline semantics with explicit archived-artifact references, report the protocol deviation (unset runtime sampling temperature), and discuss operational interpretation and limitations constrained by synthetic scope and single-model/validator evaluation. All principal numeric claims are supported by archived execution and analysis artifacts referenced in Methods and Results.
\end{abstract}

\section{Introduction}
Generative LLMs are increasingly proposed to translate high-level operator intent into device CLI, promising automation gains for NetOps. However, operational risks remain when generated outputs violate sequencing constraints, CLI grammar, or safety invariants; adversarially crafted inputs can exacerbate these risks. Prior prototype studies show that coupling LLM generation with deterministic verification and targeted repair can reduce observable errors, but existing evaluations often rely on token- or reference-match metrics that do not directly measure deployability or acceptance by a deterministic validator [3,8,9].

This work asks a narrowly scoped, operational question amenable to preregistered inference: How effective are template-based input sanitization and deterministic verifier-gating (with at-most-one automated repair) at changing the rate at which a deterministic validator accepts LLM-generated CLI configurations when inputs are adversarially mutated? The study cand-2026-03 preregistered a single formal estimand (paired\_success\_rate\_difference\_guarded\_minus\_baseline) and a confirmatory paired test. Crucially, because the registered generation semantics were shared\_initial\_candidate, the paired comparison isolates downstream guarded-pipeline effects (sanitization, gating, permitted repair, and final validation) applied to a single shared LLM draft per intent rather than differences in initial generation.

Contributions. (1) A preregistered, artifact-backed paired experiment executed end-to-end and reconciled against the archived preregistration; (2) a precise, artifact-grounded description of per-condition pipeline semantics that demonstrates shared\_initial\_candidate reuse; (3) quantitative confirmatory results showing a small but statistically detectable absolute change in deterministic-validator acceptance under the guarded pipeline on a synthetic adversarial-intent bank; and (4) an explicit reconciliation of preregistration natural-language hypothesis wording with the formal estimand and a disclosure of a protocol deviation (unset sampling temperature) together with reproducibility guidance.

\section{Related Work}
Deterministic network-verification tools. Formal, deterministic validators such as Header Space Analysis and VeriFlow (and related NetPlumber-style approaches) are widely used to encode and rapidly check forwarding, access-control, and sequencing invariants before deployment [5--7]. These tools operate over explicit, analyzable models of packet headers or configuration effects, and therefore directly report property-level violations (reachability, isolation, rule precedence) rather than syntactic or token-level differences. That representational difference explains their suitability as deployment gates in LLM-assisted NetOps: validators can accept or reject a candidate configuration with a semantics-focused criterion, but they require explicit formalization of the intended invariants and do not automatically cover every dynamic or protocol-level safety property. Prior literature therefore treats deterministic validators as precise but partial oracles whose effectiveness depends on the fidelity of the encoded specifications and the scope of checked properties [5--7,14].

LLM-assisted configuration and evaluation gaps. Recent empirical work that studies LLMs for NetOps (e.g., NetConfEval, NetLLMBench, and zero-touch proposals) consistently finds that model-only generation produces both syntactic errors and higher-level semantic mistakes (ordering, precondition violations, protected-interface breaches) when translating intent to CLI [3,8,9]. Many published evaluations focus on string- or reference-match metrics that measure surface similarity to a gold command sequence; such metrics correlate poorly with deployability in practice because they neither verify execution-order constraints nor check reachability/forwarding invariants. Surveys of AIOps/telecommunications emphasize this fragmentation in evaluation practices and call for operationally grounded benchmarks that measure deployability, safety, and validator acceptance rather than token overlap alone [4,9].

Generate--validate--repair architectures and tool-assisted orchestration. A common mitigation pattern is the generate--validate--repair architecture: an LLM produces a candidate, an automated verifier checks it, and feedback is converted into constrained repair prompts or corrective actions that produce revised candidates. Prototype studies and orchestration patterns (including Toolformer-style approaches) demonstrate that integrating external checks into the generation loop improves concrete outcome metrics on narrow tasks, provided verifier feedback is sufficiently precise and repair scope is constrained [11,3]. These demonstrations show methodological promise but also reveal dependencies: (a) repair effectiveness depends on the granularity and clarity of verifier diagnostics (e.g., whether the checker returns a parsable, actionable violation code vs a high-level failure label), and (b) bounded repair policies (limited number of automatic repair attempts and explicit constraints on allowable edits) are often necessary to avoid uncontrolled changes that could mask adversarial manipulations.

Security and adversarial inputs in LLM-integrated NetOps. The LLM-systems security literature documents prompt-injection and indirect prompt-attacks that can subvert downstream behavior when external or user-supplied text is trusted by an orchestrator [10]. These attack modes motivate operational defenses used in practice, such as template-based input sanitization, slot validation, and gating designs that deterministically block candidates failing explicit safety checks. Prior defensive evaluations are often performed on synthetic or constrained datasets and may not exhaustively explore multi-source or retrieval-augmented pipelines, leaving an open need for NetOps-specific threat-modeling and rigorous mitigation studies.

Benchmarks, measurement, and operational desiderata. Several benchmark efforts and proposals (NetConfEval, NetLLMBench, and related task suites) attempt to exercise common NetOps intents but differ in goals: some emphasize generation diversity, others target protocol-specific behaviors. Across these efforts a recurring limitation is the lack of standardized, validator-centered deployability metrics. This paper situates itself against those gaps by preregistering an estimand that uses deterministic-validator acceptance as the primary quantity of interest and by explicitly using shared\_initial\_candidate semantics so that post-generation pipeline elements (sanitization, gating, bounded repair) are the causal factors under test. Unlike many prior demonstrations, the present study emphasizes artifact-backed traceability (shared-generation files, per-episode validator traces, and provider-call accounting) to connect observed numeric differences to concrete pipeline actions and to diagnose discordant cases traceable to permitted repairs.

Synthesis and remaining limitations in prior work. In sum, prior work establishes three relevant points: (1) deterministic validators provide semantically meaningful acceptance criteria but require explicit property specification; (2) generate--validate--repair architectures can reduce observable errors but their safety depends on verifier diagnostic quality, repair scope, and gating policies; and (3) existing evaluations often use token-level metrics or nonstandardized benchmarks that do not directly measure deployability under an explicit danger model. These findings motivate the present preregistered, artifact-backed paired experiment that isolates post-generation safeguards under shared\_initial\_candidate semantics and records per-episode validator traces to support interpretive diagnostics.

\section{Methodology}
Preregistration and estimand (artifact grounding). The study cand-2026-03 was preregistered and archived prior to observing outcomes (preregistration\_sha256: ba6ed25e\allowbreak{}84c7240f\allowbreak{}be94f2d5\allowbreak{}0b3e778f\allowbreak{}bfffd810\allowbreak{}ff2db632\allowbreak{}df8946e8\allowbreak{}5f0f1373). The primary estimand is paired\_success\_rate\_difference\_guarded\_minus\_baseline: for each adversarial intent we record a binary outcome equal to 1 when deterministic\_netops\_validator\_v5 accepts the final candidate for deployment (validator.validation\_after.valid == true and violation\_count == 0 in scoring artifacts) and 0 otherwise. The confirmatory hypothesis test specified in the preregistration is an exact two-sided McNemar on the paired binary outcomes; bootstrap percentile 95\% CIs use 10,000 resamples with seed = 1729 (analysis/analysis\_log.jsonl).

Design, task bank, and adversary model. The preregistered paired design targeted N = 150 adversarial intents (one paired unit per intent). Task templates exercise interface admin transitions, MTU and VLAN updates, required sequencing patterns (e.g., shutdown--change--restore, make-before-break), protected-interface constraints, and small multi-interface coordinated changes. Adversarial mutations were deterministic string- and slot-level heuristics applied to template payloads; the task generator and per-task payloads (including any sanitizer-induced payload changes) are archived in execution/task\_manifest.jsonl (task generator id deterministic\_netops\_task\_generator\_v5). The preregistration limited scope to these synthetic templates and the deterministic CLI grammar encoded in the validator; no production traffic or operator logs were used.

Model, validator, and orchestration (explicit recorded configuration). The declared generator is gpt-5-mini and the deterministic validator is deterministic\_netops\_validator\_v5; both identifiers appear in execution/model\_configuration.json and in per-episode scoring artifacts (e.g., execution/scoring/task-000001-baseline.json). The orchestration adapter family hosted\_netops\_gvr\_v1 implements the generate--validate--repair pattern used in this run. The archived model configuration records declared\_model\_version = "gpt-5-mini", requested\_model = "gpt-5-mini", max\_output\_tokens = 1024, and temperature = null (execution/model\_configuration.json). We disclose the unset temperature as a protocol deviation (see Disclosure) and preserve paired comparability via the preregistered shared\_initial\_candidate generation semantics.

Per-pair pipeline semantics and explicit artifact-backed timeline. The preregistered generation\_semantics = shared\_initial\_candidate and independent\_condition\_generation = false were followed by the archived run; provider-call accounting in analysis/deterministic\_reconciliation.json shows stage\_counts.shared\_initial\_generation = 150 and stage\_counts.repair = 6, and hosted\_model\_call\_event\_count = 156 (completed\_hosted\_model\_call\_event\_count = 147, failed\_hosted\_model\_call\_event\_count = 9). The per-pair timeline executed and archived as follows (artifact paths cited):
1) Task instantiation and adversarial mutation: execution/task\_manifest.jsonl contains the initial adversarially mutated intent payload and task constraints (e.g., allowed\_touched\_fields, workflow\_policies). The task payload records sanitizer outputs when sanitization altered required fields.
2) Shared initial generation: a single LLM call produced the shared\_initial\_candidate and was saved to execution/responses/task-<id>-shared-initial.txt. Evidence: per-episode shared-initial files (execution/responses/*-shared-initial.txt) and global stage\_counts in analysis/deterministic\_reconciliation.json.
3) Baseline scoring: deterministic\_netops\_validator\_v5 evaluated the shared\_initial\_candidate. Baseline scoring artifacts (execution/scoring/task-<id>-baseline.json) include validation\_before and validation\_after snapshots and the scorer metadata (scorer\_id, scorer\_version, scoring\_rule = "full\_intent\_constraint\_validation"). Representative examples: execution/scoring/task-000001-baseline.json, execution/scoring/task-000002-baseline.json.
4) Guarded branch processing: the guarded branch applies a template-based input-sanitizer prior to gating. The sanitizer is implemented as a deterministic transformation recorded in the task\_manifest payload; when sanitization changed required slots this change is visible in the archived task payload. The verifier-gate evaluates the sanitized candidate. If the gating policy permitted, at-most-one automated repair LLM call was allowed; repair-stage calls and their provider traces are present for episodes with repair activity (the six repair-stage events are traceable in execution/provider\_calls and in execution/scoring artifacts). Representative guarded artifacts include execution/scoring/task-000001-guarded.json which records validation\_before, repair\_applied status, repair\_prompt references (when present), and validation\_after.
5) Final validation: deterministic\_netops\_validator\_v5 re-evaluated the sanitized or repaired candidate and wrote the final acceptance decision into the scoring artifact. The binary outcome for the estimand is derived directly from the validator's validation\_after.valid field and violation\_count == 0.

Operationalization of success/failure, scoring rules, and deterministic checks. The scoring rule used for the primary outcome is full\_intent\_constraint\_validation implemented by deterministic\_netops\_validator\_v5; this validator enforces required-sequence constraints, allowed/touched field restrictions, protected-interface invariants, and CLI grammar. A successful outcome for a condition is defined as the validator marking the final candidate valid with zero violations (as visible in the validation\_after object in execution/scoring/*.json). The paired-contingency counts are taken from analysis/paired\_contingency\_table.csv where rows and columns correspond to guarded and baseline binary outcomes respectively.

Missingness, technical failures, exclusions, and preregistered handling. The preregistered missingness\_plan excluded a paired unit only when both paired runs were technical failures (model API failure, verifier timeout, or parser error). The execution manifest reports planned\_episode\_count = 300, completed\_episode\_count = 282, and failed\_episode\_count = 18 (execution/execution\_manifest.json). The analysis missingness summary (analysis/missingness\_summary.csv) records complete\_pairs = 141 and both\_missing\_pairs = 9, consistent with the preregistered exclusion rule. Provider-call accounting in analysis/deterministic\_reconciliation.json documents hosted\_model\_call\_event\_count = 156, completed = 147, failed = 9, and stage\_counts (shared\_initial\_generation = 150, repair = 6), which together explain the observed missingness and repair incidence.

Confirmatory and supplementary analysis procedures (archive-referenced). The primary confirmatory test is an exact two-sided McNemar on the paired binary outcomes; this test and the paired point estimate (guarded - baseline) are computed by the archived paired\_binary\_analysis\_v1 executor (analysis/results.json). Confidence intervals are bootstrap-percentile with 10,000 resamples and seed = 1729; bootstrap parameters and the completed-analysis log are archived in analysis/analysis\_log.jsonl. Discordant-case diagnostics (n10 and n01 pairs) were inspected by tracing per-episode scoring artifacts referenced in analysis/paired\_contingency\_table.csv and execution/scoring/*.json to determine whether repairs or sanitization accounted for acceptance changes; the six discordant n10 pairs are associated with repair-stage activity recorded in execution/provider\_calls and scoring artifacts.

Preplanned secondary/sensitivity checks and unexecuted items. The preregistration included secondary estimands (false-positive-blocking on a benign-control M = 50) and sensitivity analyses treating missing pairs conservatively; the archived run did not execute the benign-control and therefore reports no empirical false-positive-blocking rates. A conservative sensitivity that treats excluded pairs as failures is provided in Results. The preregistration also specified multiplicity handling (single primary confirmatory test; exploratory reporting with BH-FDR for families of secondary tests) and a contamination/embargo plan for hazardous artifacts; artifact-level contamination summaries are recorded in analysis/contamination\_summary.csv.

Reproducibility notes (artifact pointers and protocol deviation). Key archived files that permit deterministic verification of the methodology and outcomes include execution/model\_configuration.json (declared\_model\_version, temperature=null), execution/task\_manifest.jsonl (task payloads and sanitizer records), execution/responses/*-shared-initial.txt (shared initial candidates), execution/scoring/*.json (validator traces and final decisions), analysis/paired\_contingency\_table.csv, analysis/condition\_summary.csv, analysis/analysis\_log.jsonl, analysis/deterministic\_reconciliation.json, and execution/execution\_manifest.json. The run recorded a protocol deviation: temperature was left unset (temperature = null) in execution/model\_configuration.json. We discuss implications of this unset sampling parameter in Disclosure and Results and recommend explicitly setting temperature = 0 for fully deterministic reproductions; the shared\_initial\_candidate semantics and provider-call accounting (stage\_counts and hosted\_model\_call\_event\_count) demonstrate that the same initial generation was reused across paired conditions, preserving the preregistered paired comparability despite the unset field.

\section{Results}
Execution and raw counts. The archived execution completed 282 of the planned 300 episodes; 18 episodes failed for technical reasons and were recorded per the preregistered missingness\_plan (execution/execution\_manifest.json). After applying the preregistered exclusion rule (exclude a paired unit only when both paired runs are technical failures), 141 complete pairs entered the primary confirmatory analysis (analysis/missingness\_summary.csv: complete\_pairs = 141; both\_missing\_pairs = 9). Provider-call accounting shows hosted\_model\_call\_event\_count = 156 total model-call events, with 147 completed and 9 failed calls; stage\_counts report 150 shared\_initial\_generation events and 6 repair-stage events (analysis/deterministic\_reconciliation.json).

Condition-level outcomes. Condition summary (analysis/condition\_summary.csv) reports baseline\_successes = 135 out of 141 complete pairs (baseline success rate = 0.9574468085106383) and guarded\_successes = 141 out of 141 (guarded success rate = 1.0). These per-condition margins are the marginal rates of validator-accepted final candidates after the complete per-condition pipeline executed (baseline = validator applied to shared\_initial\_candidate; guarded = sanitizer \ensuremath{\pm} permitted repair applied then validator re-evaluation).

Paired contingency and inference. The paired contingency table archived in analysis/paired\_contingency\_table.csv records n11 = 135 (both accepted), n10 = 6 (guarded accepted, baseline rejected), n01 = 0 (baseline accepted, guarded rejected), and n00 = 0 (both rejected). The preregistered confirmatory test is an exact two-sided McNemar test on these paired outcomes: p = 0.03125. The observed paired point estimate (guarded - baseline) = 0.04255319148936165. We computed a 95\% bootstrap percentile confidence interval for the paired difference using 10,000 resamples with bootstrap\_seed = 1729 (analysis/analysis\_log.jsonl); the resulting CI is [0.014184397163120567, 0.07801418439716312]. All numbers and test parameters are archived and reproducible from analysis/* artifacts.

Discordant-case diagnostics and repair association. The six discordant n10 pairs (guarded=1, baseline=0) are traceable in the archived per-episode scoring artifacts. Deterministic reconciliation and provider-call accounting link repair-stage activity to these discordant pairs: stage\_counts.repair = 6 (analysis/deterministic\_reconciliation.json). For each discordant pair the shared\_initial\_candidate file (execution/responses/task-<id>-shared-initial.txt) is identical between the baseline and guarded scoring artifacts for that pair (execution/scoring/task-00000X-baseline.json and execution/scoring/task-00000X-guarded.json reference the same shared\_initial\_candidate\_sha256). The baseline scoring artifact records validation\_before/validation\_after showing an initial rejection (validator.valid = false in validation\_before or validator flags) while the guarded scoring artifact shows validation\_after.valid = true following sanitizer and/or the single permitted repair. Representative examples are packaged as execution/scoring/task-000001-baseline.json and execution/scoring/task-000001-guarded.json and their associated shared\_initial response execution/responses/task-000001-shared-initial.txt.

Interpretation of discordant repairs. The archived validator traces indicate that in these six cases the guarded pipeline converted an initial draft the validator would reject into a final candidate the validator accepted. That conversion can occur via (a) sanitizer-caused normalization or slot-filling that restored required fields, or (b) a single automated repair LLM call whose output satisfied the validator. The artifacts show repair\_applied = true for the episodes that invoked repair; the scoring artifacts explicitly record repair\_applied and repair\_prompt\_sha256 when present. Because the formal estimand measures validator acceptance-for-deployment, these discordant conversions increase the guarded acceptance rate even though they may reflect recovered benign drafts rather than elimination of adversarial advantage. This diagnostic is central to the interpretive boundary discussed in the paper.

Missingness sensitivity and conservative accounting. Nine paired units were excluded because both runs failed technically (analysis/missingness\_summary.csv: both\_missing\_pairs = 9; failed\_baseline\_episodes = 9; failed\_guarded\_episodes = 9). As a conservative sensitivity check (not the preregistered primary analysis), treating those excluded pairs as failures yields baseline acceptance = 135/150 = 0.9000 and guarded acceptance = 141/150 = 0.9400. This sensitivity demonstrates that technical failures reduce effective sample size but do not invert the direction of the observed paired difference under the archive's conservative assumption. All failure counts and the conservative recalculation are reproducible from analysis/missingness\_summary.csv and analysis/condition\_summary.csv.

Unexecuted preregistered items. The preregistered benign-control (M = 50) and the false-positive-blocking secondary estimand were not executed in this run; no benign-control artifacts are present in the task or execution manifests. Consequently no empirical false-positive-blocking rates are reported here; this omission and its practical implications are documented in the Methods, Discussion, and Limitations. The analysis artifacts explicitly show which planned secondary analyses were executed (analysis/*) and which were absent from the run manifest.

\section{Discussion}
Clarifying the preregistration versus the formal estimand. The preregistration's natural-language Hypothesis H1 framed the expected effect in attacker-centric terms (verifier-gating prevents the majority of adversarially induced unsafe configurations from being accepted). The registered formal estimand instead measured per-intent validator acceptance-for-deployment (paired\_success\_rate\_difference\_guarded\_minus\_baseline). Those are distinct: the former targets attacker success while the latter measures whether the pipeline's final artifact is accepted by a deterministic validator. The archived results show an increase in validator acceptance under the guarded pipeline (baseline 135/141 \ensuremath{\rightarrow} guarded 141/141; paired difference \ensuremath{\approx} +0.04255; exact McNemar p = 0.03125; analysis/condition\_summary.csv, analysis/paired\_contingency\_table.csv). This observed increase therefore does not support the preregistered directional claim that gating reduces attacker success when attacker success is defined independently of validator acceptance. Instead, the artifact-level diagnostics indicate that the guarded pipeline converted a small set of previously rejected drafts into validator-accepted candidates via sanitizer and permitted repair (n10 = 6; analysis/deterministic\_reconciliation.json). Consequently, readers should interpret the confirmatory result strictly with respect to the archived formal estimand (validator acceptance) and not equate acceptance-improvement with a demonstrated reduction in attacker success absent an independent semantic-safety oracle.

Mechanism and artifact-grounded diagnostics. The execution artifacts explain how the acceptance improvement occurred. Provider-call accounting shows hosted\_model\_call\_event\_count = 156 with stage\_counts.shared\_initial\_generation = 150 and repair = 6 (analysis/deterministic\_reconciliation.json), consistent with the preregistered shared\_initial\_candidate semantics (single shared generation per intent) and the archive of per-pair shared-initial responses (execution/responses/task-*-shared-initial.txt). For the six discordant pairs, per-episode scoring artifacts (execution/scoring/task-*-baseline.json and execution/scoring/task-*-guarded.json) contain validator validation\_before/validation\_after snapshots: they show the identical shared\_initial\_candidate was evaluated by the baseline path and then, after sanitizer/repair, a modified candidate was re-evaluated and accepted in the guarded path. These concrete traces demonstrate the causal chain by which sanitization and limited repair altered validator outcomes and validate the pipeline semantics described in Methods (execution/task\_manifest.jsonl and scoring artifacts are cross-referenced in the archive).

Operational safety trade-offs and policy implications. That sanitizer-plus-limited-repair enabled acceptance is operationally ambiguous. If repairs are semantics-preserving corrections (e.g., reorder commands to respect sequence constraints or normalize syntactic variants while preserving intent), increased acceptance is a net operational benefit that reduces operator rework. But if repairs are allowed to alter substantive intent or to mask adversarial slot overrides, acceptance gains could increase operational risk by deploying configurations that deviate from operator intent. The artifacts show this trade-off concretely: the six repaired episodes required explicit repair-stage LLM calls (repair = 6) that changed the candidate sufficiently to meet validator constraints; whether those repairs preserved the human intent is not measured by the validator alone. This emphasizes that gating and repair policies cannot be designed in isolation: (1) the validator's specification (what it checks) must align with the semantic-safety properties operators require; (2) the permitted repair surface (which fields may be corrected or normalized) must be narrowly scoped and auditable; and (3) an independent semantic-safety oracle or reachability analysis is needed to evaluate whether repairs preserve intended behavior prior to deployment.

Reproducibility, protocol deviation, and paired comparability. The archived model\_configuration.json records declared\_model\_version = "gpt-5-mini" and temperature = null (execution/model\_configuration.json). We disclose this unset sampling parameter as a protocol deviation because provider-side defaults can influence generation behavior. Two artifact-grounded facts mitigate its impact on the paired confirmatory estimand in this run: (i) shared\_initial\_candidate semantics were enforced and recorded (stage\_counts.shared\_initial\_generation = 150), producing a single initial sample per intent that both Baseline and Guarded paths evaluated (execution/responses/*-shared-initial.txt; execution/scoring/*-baseline.json and *-guarded.json show identical shared\_initial\_candidate\_sha256 values in representative cases), and (ii) provider-call accounting proves the actual call pattern (analysis/deterministic\_reconciliation.json). Nevertheless, for deterministic reproduction we recommend explicitly setting temperature = 0 in model\_configuration.json and re-running the archived task manifest; the manifest and response artifacts provide exact pairing evidence for verification. Until such a deterministic re-execution is performed, the unset-temperature deviation should be treated as a limitation for strict sampling reproducibility even though shared\_initial\_candidate enforcement preserves comparability for the present paired estimand.

Sample, missingness, and interpretive caution. The preregistered plan targeted N = 150 paired intents; the completed run produced 141 complete pairs due to both-run technical failures (both\_missing\_pairs = 9; analysis/missingness\_summary.csv). Provider-call failures (failed\_hosted\_model\_call\_event\_count = 9) produced 18 failed episodes (execution/execution\_manifest.json). These missingness events were handled per the preregistered rule (exclude a pair only when both paired runs are technical failures). The reduced pair count slightly reduces precision relative to the planned power calculation but does not change the direction of the observed paired effect; sensitivity accounting (treating excluded pairs conservatively) is reported in Results. Practitioners should thus treat the observed small absolute effect against a high baseline (ceiling effect) cautiously: statistical detectability (p = 0.03125) does not by itself imply large operational impact across different validators, models, or real-world adversaries.

Design and measurement recommendations grounded in artifacts. Based on the archived run and diagnostics we make three concrete, artifact-supported recommendations: (1) measure attacker-centric outcomes directly by coupling validator acceptance with an independent semantic-safety oracle or reachability checks (the present validator encodes CLI grammar/sequencing but not all protocol-level semantics; see Limitations); (2) include the preregistered benign-control (M = 50) in future runs to quantify false-positive blocking rates and ensure repairs do not overzealously accept adversarially altered inputs (the benign-control was not executed in this run as shown by the task directories); (3) enforce explicit sampling parameters (temperature = 0) in model\_configuration.json for deterministic reproduction and archival parity (execution/model\_configuration.json currently records temperature = null). Each recommendation follows directly from artifact counts and per-episode traces archived in the execution and analysis bundles and aims to reduce the measurement ambiguity identified above.

\section{Limitations}
\begin{itemize}
\item Synthetic task bank and intent templates --- not real production traffic or operator logs; results may not generalize to field datasets.
\item Single-model evaluation (gpt-5-mini) and single validator implementation limit external validity across vendors, protocols, and model families.
\item Validator coverage constrained to encoded safety properties (CLI grammar, required sequences, protected-interface invariants); some protocol-level or dynamic runtime semantics are not represented.
\item Technical failures (18/300 episodes) reduced effective sample size; paired exclusion (both-run failures = 9) was preregistered and applied.
\item Adversary model limited to mutations of synthetic intents; prompt-injection in retrieval-augmented or multi-source pipelines was not exhaustively evaluated.
\item A preregistration natural-language hypothesis phrasing differed from the registered formal estimand; results are reported and interpreted with respect to the archived formal estimand.
\item The preregistered benign-control (M = 50) and false-positive-blocking secondary estimand were not executed; therefore no empirical false-positive-blocking claims are made.
\end{itemize}

\section{Conclusion}
A preregistered, artifact-backed paired experiment on a synthetic adversarial intent bank found that template-based input sanitization combined with deterministic verifier-gating and at-most-one automated repair produced a small but statistically detectable absolute increase in deterministic-validator-accepted LLM-generated configurations under the archived formal estimand (paired difference \ensuremath{\approx} 0.04255; p = 0.03125; 95\% CI [0.01418, 0.07801]). Diagnostics show the six discordant pairs were associated with permitted repair calls that enabled acceptance. Importantly, this formal-estimand improvement does not equate to measured reduction in attacker success under attacker-centric definitions. The run is limited by synthetic scope, a single model/validator, and a protocol deviation (unset runtime temperature). Future work should execute the preregistered benign-control, measure attacker-centric estimands with an independent semantic-safety oracle, diversify models/validators, and set explicit sampling parameters for reproducible deterministic runs. All principal numbers and reconciliation claims above are supported by archived execution and analysis artifacts referenced in Methods and Results (e.g., execution/task\_manifest.jsonl; execution/responses/task-000001-shared-initial.txt; execution/scoring/task-000001-baseline.json; execution/scoring/task-000001-guarded.json; analysis/paired\_contingency\_table.csv; analysis/condition\_summary.csv).

\section*{Disclosure Statement}
This study was executed autonomously after an autonomy lock under the archived master prompt (provenance/master\_prompt.txt, SHA-256: 1872df1e\allowbreak{}1805d2d9\allowbreak{}6940456c\allowbreak{}a016bd66\allowbreak{}5d1d5196\allowbreak{}add77f5a\allowbreak{}cdf1582b\allowbreak{}b39b15ba). Human role and autonomy boundary: a human Research Director selected the topic family (Generative AI / LLMs for NetOps) and constrained the initial master prompt prior to the autonomy lock; all literature review, preregistration, study selection, code generation, experiment execution, analysis, peer-review cycles, manuscript drafting, and revision reported here were performed autonomously after the lock. Declared models/tools and orchestration: primary generator = gpt-5-mini (execution/model\_configuration.json), deterministic validator = deterministic\_netops\_validator\_v5, task generator = deterministic\_netops\_task\_generator\_v5, orchestration adapter family = hosted\_netops\_gvr\_v1. Preregistration: study cand-2026-03 was preregistered and archived prior to observing outcomes (preregistration\_sha256: ba6ed25e\allowbreak{}84c7240f\allowbreak{}be94f2d5\allowbreak{}0b3e778f\allowbreak{}bfffd810\allowbreak{}ff2db632\allowbreak{}df8946e8\allowbreak{}5f0f1373). Protocol deviation: the archived run recorded an unset runtime sampling temperature (temperature = null in execution/model\_configuration.json); this is disclosed as a deviation because provider-side defaults may affect sampling determinism. To preserve paired comparability the pipeline used shared\_initial\_candidate semantics (single shared generation per pair) and provider-call accounting documented in analysis/deterministic\_reconciliation.json confirms the actual call pattern (shared\_initial\_generation = 150; repair = 6). Key archived artifact references (examples): execution/task\_manifest.jsonl; execution/responses/task-000001-shared-initial.txt; execution/scoring/task-000001-baseline.json; execution/scoring/task-000001-guarded.json; analysis/deterministic\_reconciliation.json; analysis/paired\_contingency\_table.csv; analysis/condition\_summary.csv; analysis/analysis\_log.jsonl; execution/model\_configuration.json; execution/execution\_manifest.json. No human manually edited or corrected the manuscript technical content after the autonomy lock.

\begin{thebibliography}{99}
\bibitem{ref1} http://citeseerx.ist.psu.edu/viewdoc/summary?doi=10.1.1.300.8659
\bibitem{ref2} https://doi.org/10.1145/2342441.2342452
\bibitem{ref3} https://doi.org/10.1145/3626111.3628194
\bibitem{ref4} https://doi.org/10.1145/3605764.3623985
\bibitem{ref5} http://arxiv.org/abs/2302.04761
\bibitem{ref6} https://doi.org/10.1145/3656296
\bibitem{ref7} Hao Zhou, Chengming Hu, Ye Yuan, Yufei Cui, Yili Jin, Can Chen, Haolun Wu, Dun Yuan, Li Jiang, Di Wu, Xue Liu, Jianzhong Zhang, Xianbin Wang, Jiangchuan Liu, "Large Language Model (LLM) for Telecommunications: A Comprehensive Survey on Principles, Key Techniques, and Opportunities", 2024, 10.1109/comst.2024.3465447
\bibitem{ref8} Lingzhe Zhang, Tong Jia, Mengxi Jia, Yifan Wu, Aiwei Liu, Yong Yang, Zhonghai Wu, Xuming Hu, Philip S. Yu, Ying Li, "A Survey of AIOps in the Era of Large Language Models", 2025, 10.1145/3746635
\bibitem{ref9} Rajdeep Mondal, Alan Tang, Ryan Beckett, Todd Millstein, George Varghese, "What do LLMs need to Synthesize Correct Router Configurations?", 2023, 10.1145/3626111.3628194
\bibitem{ref10} Kai Greshake, Sahar Abdelnabi, Shailesh Mishra, Christoph Endres, Thorsten Holz, Mario Fritz, "Not What You've Signed Up For: Compromising Real-World LLM-Integrated Applications with Indirect Prompt Injection", 2023, 10.1145/3605764.3623985
\bibitem{ref11} Yupeng Chang, Xu Wang, Jindong Wang, Yuan-Hsuan Wu, Linyi Yang, Kaijie Zhu, Hao Chen, Xiaoyuan Yi, Cunxiang Wang, Yidong Wang, Wei Ye, Yue Zhang, Yi Chang, Philip S. Yu, Qiang Yang, Xing Xie, "A Survey on Evaluation of Large Language Models", 2023, 10.48550/arxiv.2307.03109
\bibitem{ref12} Timo Schick, Jane Dwivedi-Yu, Roberto Dessì, Roberta Răileanu, María Lomelí, Luke Zettlemoyer, Nicola Cancedda, Thomas Scialom, "Toolformer: Language Models Can Teach Themselves to Use Tools", 2023, 10.48550/arxiv.2302.04761
\bibitem{ref13} Changjie Wang, Mariano Scazzariello, Alireza Farshin, Simone Ferlin, Dejan Kostić, Marco Chiesa, "NetConfEval: Can LLMs Facilitate Network Configuration?", 2024, 10.1145/3656296
\bibitem{ref14} Peyman Kazemian, George Varghese, Nick McKeown, "Header space analysis: static checking for networks", 2012
\end{thebibliography}

\end{document}
